\section{Conclusion}
\label{sec:conclusion}

Using a contextual skill and reusing its internal states are distinct
capabilities. Across four models and two benchmarks, state transfer
weakens before the original execution's dependence on skill attention,
with half-effect boundaries separated by $8$ to $20$ layers.
Early content controls establish that the transferred states carry
skill-relevant information, and the depth comparisons show that this
information becomes harder to reuse while the skill remains useful.

Source-cache replacement provides a way to recover part of that transfer.
At layer $16$, adding caches improves accuracy over states alone by
$6.8$ to $43.4$ percentage points in all eight comparisons, while cutting
subsequent skill reading reduces the joint-repair accuracy in each.
Complete-source controls further show that skill-position KV alone is
insufficient to recover donor performance. Full-prefix copying reproduces
the donor, while the depth effect of joint repair varies across models
and tasks.
The mathematical analysis connects cache alignment to output stability
under explicit conditions. Together, these findings motivate evaluating
state reuse separately from continued skill use and including cached
information in the design of reusable skill representations.
